% RESEARCH_DRAFT_SHA256: 45ec5d14fdf622015d4561ef03d721db631fbddf00bc6c2d943ff48c779c3317
% Layout and source-reference wrappers expanded for the web renderer.
\documentclass[11pt]{article}
\usepackage[margin=25mm]{geometry}
\usepackage{amsthm}
\usepackage{fontspec,unicode-math}
\setmainfont{Latin Modern Roman}
\setmathfont{Latin Modern Math}
\usepackage{amsmath,amssymb,mathtools}
\usepackage{xurl,hyperref}
\hypersetup{colorlinks=true,urlcolor=blue,linkcolor=blue}
\setcounter{secnumdepth}{0}
\setlength{\parindent}{0pt}
\setlength{\parskip}{5pt}
\setlength{\emergencystretch}{3em}
\newcommand{\sref}[2]{\hyperlink{source:#1:#2}{[S#2]}}
\newcommand{\cataloguescope}[1]{\paragraph{Source question.}#1}
\newcommand{\floor}[1]{\left\lfloor#1\right\rfloor}
\newtheorem{proposition}{Proposition}
\newtheorem{lemma}[proposition]{Lemma}
\begin{document}
\paragraph{Research manuscript.} OpenAI \texttt{gpt-6-astra}, reasoning \texttt{ultra}; 11 October 2026.
\paragraph{Preserved source body.} The complete catalogue mathematical text and sources follow with layout and source links expanded, including source qualifications and its historical preparation notes. New research begins at the section headed Research attempt.

\section{Circuit-Size Hierarchy in Bounded Arithmetic}\label{statement}

\subsection{Definitions and mathematical statement}

\paragraph{Shared notation.}Standard number systems are denoted by $\mathbb N,\mathbb Z,\mathbb Q,\mathbb R,\mathbb C$; index ranges, computational models and specialized notation follow the cited problem source and the definitions below.


A Boolean circuit is an acyclic graph of Boolean gates computing a function on $n$ input bits; circuit size counts gates. A circuit-size hierarchy theorem separates classes of functions computable within different asymptotic circuit-size budgets, under the source gate and encoding conventions. The bounded-arithmetic theory $\mathrm{PV}_1$ formalizes polynomial-time reasoning with symbols for polynomial-time functions and its specified induction axioms. A language is in polynomial time if membership is decided by a uniform polynomial-time algorithm. An explicit circuit lower bound supplies such a uniform language together with a proved lower bound for its circuit sizes.
\paragraph{Statement recorded in the supplied source.}
Does the bounded-arithmetic theory $\mathrm{PV}_1$ prove the Circuit-Size Hierarchy theorem? If it does, can one extract from such a proof explicit polynomial-time languages with the corresponding circuit lower bounds? \href{https://doi.org/10.4171/owr/2024/15}{[S2]}
\paragraph{Scope qualification.}
The formal hierarchy statement, size bounds and proof-extraction claim follow the cited source. A nonconstructive existence proof does not by itself produce a uniform polynomial-time hard language.
\paragraph{Source status.}
The supplied catalogue labels this record open. Its source-status metadata was checked against the supplied snapshot on 2026-10-11. This import does not independently establish current open status or mathematical correctness. The catalogue literature assessment is dated 2026-08-22.
\subsection{Short English statement}
Can polynomial-time bounded arithmetic prove a circuit-size hierarchy and yield explicit efficiently decidable functions realizing its lower bounds?
\subsection{Sources}
\begin{itemize}
\item[S1] Ulam AI and the UnsolvedMath Contributors, UnsolvedMath: A Curated Collection of Open Mathematics Problems, catalogue entry. Catalogue curation and metadata: CC BY 4.0. The supplied source question is retained; definitions and the English summary are editorial additions.
\url{https://huggingface.co/datasets/ulamai/UnsolvedMath}.
\item[S2] Proof Complexity and Beyond (2024), Oberwolfach Reports, DOI 10.4171/owr/2024/15, pp. 871--934
\url{https://doi.org/10.4171/owr/2024/15}.
\item[S3] Carmosino et al., Provability of the Circuit Size Hierarchy
\url{https://doi.org/10.4230/LIPIcs.ITCS.2025.30}.
\item[S4] Oberwolfach Report 15/2024
\url{https://oa.tib.eu/renate/bitstreams/b0ef46dc-1bdf-45b7-a4f5-ca0a9e60a7d6/download}.
\textit{Underlying Oberwolfach Reports material retains its source-specific licence, including CC BY-SA 4.0 where applicable; see THIRD\_PARTY\_NOTICES.md.}
\end{itemize}
\paragraph{Preparation and review.}
Original question attributed to the cited source and UnsolvedMath contributors. Codex (GPT-6; exact serving revision not exposed) assisted with definitions, English summaries, source inspection and formatting. No new human mathematical authorship or independent proof verification is asserted. Catalogue identity checked against the complete supplied registry and existing repository statements; known flagged repetitions and obvious identical conjectures excluded. Full mathematical-equivalence/literature review has not been performed. Statement preparation and final review followed contribution prompts C and D; automated checks verify structure and identity, not open status.
\label{end:statement}

\section{Research attempt}
\subsection{Literature, source context and scope}
The full original catalogue body and source qualifications are preserved above. The relevant primary workshop passages were read in the cached publisher report. Historical claims and subsequent attributed results are distinguished from the mathematical work below. A restricted example, standard lemma or known theorem is not counted as new progress on the entire original target. Missing assumptions and repeated or rhetorical extracts are not silently turned into different solved problems.

\subsection{A classical hierarchy argument and its formalization gap}
The source asks for provability inside weak arithmetic, not merely a classical existence proof. The 2025 primary paper reports circuit hierarchy results in a stronger theory and a restricted formula hierarchy in $\mathrm{PV}_1$; those do not establish the full recorded $\mathrm{PV}_1$ claim. We attempt a direct counting-and-interpolation route and identify the non-effective step.
Count fan-in-two Boolean circuits with at most $s$ gates, allowing all 16 binary truth tables, two constants, and an output chosen among inputs, constants and gates. A topological encoding gives at most
\[
(s+1)(n+s+2)[16(n+s+2)^2]^s
\]
circuits: sum the bounds for each number of gates from zero to $s$, overcounting predecessor choices by $n+s+2$.
\begin{proposition}
For $n\ge10$, some $n$-input Boolean function has circuit complexity exceeding $s=\lfloor2^n/(10n)\rfloor$. Consequently, for fixed real $a>b\ge1$ and all sufficiently large $n$, there exists a Boolean function of circuit complexity strictly between $n^b$ and $n^a$.
\end{proposition}
\begin{proof}
For the stated $s$, both $s+1$ and $n+s+2$ are at most $2^n$. Hence the base-two logarithm of the count is at most $2n+s(2n+4)$. For $n\ge10$, $2n\le2^n/10$ (true at ten and preserved by increasing $n$), and $s(2n+4)\le(6/25)2^n$. The count is therefore less than $2^{2^n}$, the number of truth tables. This proves the first assertion.
For the second, put $S=n^b$. Eventually $S<s$. Choose a function $f$ of complexity greater than $s$, list its accepting inputs, and form partial disjunctions of the corresponding exact-input minterms, starting at the constant-zero function. Adding a minterm requires at most $n$ negations, $n-1$ conjunctions and one disjunction, at most $2n$ new gates. Thus the first partial function $h$ whose minimum circuit complexity exceeds $S$ has complexity at most $S+2n$: the preceding partial function has a circuit of size at most $S$ to which the new minterm can be added. Since $a>b\ge1$, $n^b+2n<n^a$ for all sufficiently large $n$. This $h$ has the required intermediate size.
\end{proof}
The definition of the first crossing uses the minimum circuit complexity of successive exponentially described truth tables. The proof supplies no polynomial-time method for testing that condition, no polynomially bounded sequence of formal witnesses, and no derivation in $\mathrm{PV}_1$. Replacing this semantic search by an efficient construction is the first unsupported inference in an attempted weak-arithmetic proof. Classical counting alone does not license it.
\subsection{Verification and remaining gaps}
The numerical inequalities were checked over a range of $n$ using integer/rational bounds, and small truth tables verify that adding one minterm changes the represented function exactly as stated. These are ordinary finite computations outside a weak-theory proof. The proposition is a classical auxiliary existence argument, not a new circuit lower bound or a resolution of the requested formal provability.

\section{References and source attribution}
Primary workshop report: \url{https://doi.org/10.4171/owr/2024/15}. Relevant passages read: 2024/15 pp.928--929, circuit-size hierarchy and its weak-arithmetic formulation.

\url{https://drops.dagstuhl.de/entities/document/10.4230/LIPIcs.ITCS.2025.30}. Read primary abstract on circuit and formula size hierarchies in bounded arithmetic; stronger-theory and restricted-formula results are not identified with the full PV1 question.

All retrievals were checked on 11 October 2026. Original question attribution remains in the preserved Sources section. External theorems are attributed and not claimed as new discoveries or fully verified merely by citation.


\paragraph{Integration review.} Codex (GPT-6; exact serving revision not exposed) reviewed the written arguments and expanded layout and source-reference wrappers. Repository metadata and the production web renderer were checked using Python 3.13 and JavaScript with MathJax 3.2.2. This model-assisted review is not independent human proof verification; no PDF compilation or formal proof checking is asserted.

\section{Completion Estimate}
\noindent\textbf{COMPLETION: 0\%}
No substantiated new progress toward the entire original target is claimed. The worked results and failed route are delimited above; a known result, illustrative model, or auxiliary obstruction is not counted as original-target completion.

\section{Final status and attribution}
\textbf{Attempt status: unresolved.} This describes this research attempt, not a new ruling on the source catalogue's status or on external literature claims.
The mathematics and this self-audit were generated by OpenAI \texttt{gpt-6-astra}, reasoning \texttt{ultra}, on 11 October 2026. Codex (GPT-6; exact serving revision not exposed) orchestrated the batch, read and reviewed its written mathematical arguments, and contributed formatting/repository review. No human mathematical authorship, independent human proof checking, or formal verification is claimed. The historical preparation paragraphs in the reproduced catalogue body concern that body only.

Contribution prompt A was applied to the complete definition and supplied source record, followed by the adversarial checks of prompt B. Local reproduction files are \texttt{batch025/checks.py}, \texttt{computations.json}, and the per-ID audit JSON. They were not public links at preparation time. Numerical checks are finite and do not prove unrestricted results. No LaTeX engine was found locally; compilation was not performed by this research agent.

\end{document}
